\chapter{Combinatoire et Dénombrement}

Ce chapitre traite du dénombrement, c'est-à-dire de l'art de compter le nombre d'éléments d'un ensemble fini sans pour autant les énumérer tous. Les outils développés ici (permutations, arrangements, combinaisons) sont fondamentaux pour le calcul des probabilités.

\section{Principes Fondamentaux du Dénombrement}

Soit $E$ un ensemble fini. On note $\text{Card}(E)$ (ou $\#E$) le cardinal de $E$, c'est-à-dire le nombre d'éléments de $E$.

\subsection{Principe additif}

Le principe additif concerne la réunion d'ensembles disjoints.

\begin{theoreme}[Principe additif]
Si $A$ et $B$ sont deux sous-ensembles d'un ensemble $E$, disjoints ($A \cap B = \emptyset$), alors :
\[
\text{Card}(A \cup B) = \text{Card}(A) + \text{Card}(B)
\]
De façon générale, si $A$ et $B$ sont quelconques :
\[
\text{Card}(A \cup B) = \text{Card}(A) + \text{Card}(B) - \text{Card}(A \cap B)
\]
\end{theoreme}

\subsection{Principe multiplicatif}

Le principe multiplicatif s'applique lorsqu'on réalise une succession de choix indépendants.

\begin{theoreme}[Principe multiplicatif]
Soit un ensemble de choix successifs où le premier choix offre $n_1$ possibilités, le deuxième $n_2$ possibilités, ..., et le $k$-ième choix offre $n_k$ possibilités.
Le nombre total de configurations possibles est :
\[
N = n_1 \times n_2 \times \dots \times n_k
\]
\end{theoreme}

\begin{definition}[$p$-uplet ou liste]
Soit $E$ un ensemble à $n$ éléments. Un **$p$-uplet** (ou liste de longueur $p$) d'éléments de $E$ est une suite ordonnée de $p$ éléments de $E$.
Le nombre de $p$-uplets d'un ensemble à $n$ éléments est :
\[
n^p
\]
\end{definition}

\section{Arrangements et Permutations}

\subsection{Arrangements}

Un arrangement est un choix ordonné d'éléments distincts.

\begin{definition}[Arrangement]
Soit $E$ un ensemble à $n$ éléments. Un **arrangement de $p$ éléments de $E$** ($0 \le p \le n$) est un $p$-uplet d'éléments **distincts** de $E$.
Le nombre d'arrangements de $p$ éléments parmi $n$ est noté $A_n^p$ et vaut :
\[
A_n^p = n \times (n-1) \times \dots \times (n-p+1) = \frac{n!}{(n-p)!}
\]
\end{definition}

\subsection{Permutations}

\begin{definition}[Permutation et factorielle]
Soit $E$ un ensemble à $n$ éléments. Une **permutation** de $E$ est un arrangement de $n$ éléments de $E$ (c'est-à-dire un ordre sur tous les éléments de $E$).
Le nombre de permutations d'un ensemble à $n$ éléments est noté $n!$ (lire « factorielle $n$ ») et vaut :
\[
n! = n \times (n-1) \times \dots \times 2 \times 1
\]
Par convention, $0! = 1$.
\end{definition}

\section{Combinaisons}

Contrairement aux arrangements, l'ordre n'importe pas dans une combinaison (comme pour une main de cartes).

\subsection{Définition}

\begin{definition}[Combinaison]
Soit $E$ un ensemble à $n$ éléments. Une **combinaison de $p$ éléments de $E$** ($0 \le p \le n$) est un sous-ensemble de $E$ contenant $p$ éléments (sans ordre, sans répétition).
Le nombre de combinaisons de $p$ éléments parmi $n$ est noté $\binom{n}{p}$ (lire « $p$ parmi $n$ ») et vaut :
\[
\binom{n}{p} = \frac{A_n^p}{p!} = \frac{n!}{p!(n-p)!}
\]
\end{definition}

\subsection{Propriétés des coefficients binomiaux}

\begin{propriete}[Relations de symétrie et Pascal]
Pour tous entiers $n$ et $p$ tels que $0 \le p \le n$ :
\begin{itemize}
    \item \textbf{Symétrie} : $\binom{n}{p} = \binom{n}{n-p}$
    \item \textbf{Cas particuliers} : $\binom{n}{0} = 1$, $\binom{n}{1} = n$, $\binom{n}{n} = 1$
    \item \textbf{Formule de Pascal} (pour $0 \le p \le n-1$) :
    \[
    \binom{n}{p} + \binom{n}{p+1} = \binom{n+1}{p+1}
    \]
\end{itemize}
\end{propriete}

\subsection{Le Triangle de Pascal}

La formule de Pascal permet de calculer de proche en proche les coefficients binomiaux sous forme triangulaire :

\begin{center}
\begin{tikzpicture}[scale=1.0, every node/.style={circle, draw=defcolor!50, fill=defbg, minimum size=24pt, inner sep=0pt, font=\sffamily\small}]
  % Row 0
  \node (n00) at (0, 0) {1};
  
  % Row 1
  \node (n10) at (-0.6, -0.8) {1};
  \node (n11) at (0.6, -0.8) {1};
  
  % Row 2
  \node (n20) at (-1.2, -1.6) {1};
  \node (n21) at (0, -1.6) {2};
  \node (n22) at (1.2, -1.6) {1};
  
  % Row 3
  \node (n30) at (-1.8, -2.4) {1};
  \node (n31) at (-0.6, -2.4) {3};
  \node (n32) at (0.6, -2.4) {3};
  \node (n33) at (1.8, -2.4) {1};
  
  % Row 4
  \node (n40) at (-2.4, -3.2) {1};
  \node (n41) at (-1.2, -3.2) {4};
  \node (n42) at (0, -3.2) {6};
  \node (n43) at (1.2, -3.2) {4};
  \node (n44) at (2.4, -3.2) {1};
  
  % Fleches d'addition de Pascal
  \draw[->, >=latex, thin, propcolor] (n20) -- (n31);
  \draw[->, >=latex, thin, propcolor] (n21) -- (n31);
  \draw[->, >=latex, thin, propcolor] (n21) -- (n32);
  \draw[->, >=latex, thin, propcolor] (n22) -- (n32);
  
  \draw[->, >=latex, thin, propcolor] (n31) -- (n42);
  \draw[->, >=latex, thin, propcolor] (n32) -- (n42);
  
  % Labels
  \node[draw=none, fill=none, left=30pt] at (n00) {\footnotesize $n=0$};
  \node[draw=none, fill=none, left=30pt] at (n10) {\footnotesize $n=1$};
  \node[draw=none, fill=none, left=30pt] at (n20) {\footnotesize $n=2$};
  \node[draw=none, fill=none, left=30pt] at (n30) {\footnotesize $n=3$};
  \node[draw=none, fill=none, left=30pt] at (n40) {\footnotesize $n=4$};
\end{tikzpicture}
\end{center}

\section{Méthodes Clés}

\begin{methode}[Choisir le bon outil de dénombrement]
Pour savoir quelle formule utiliser pour dénombrer les tirages ou configurations :
\begin{enumerate}
    \item **L'ordre compte-t-il ?**
    \begin{itemize}
        \item **Non** : Il s'agit de **combinaisons** $\implies \binom{n}{p}$.
        \item **Oui** : Passer à la question 2.
    \end{itemize}
    \item **Y a-t-il des répétitions ?**
    \begin{itemize}
        \item **Oui** : Il s'agit de **$p$-uplets (listes)** $\implies n^p$.
        \item **Non** : Il s'agit d'un **arrangement** ($A_n^p$) ou d'une **permutation** ($n!$ si $p=n$).
    \end{itemize}
\end{enumerate}
\end{methode}

\begin{exemple}[Application des méthodes]
Un digicode est composé d'une lettre parmi A, B, C suivie de 4 chiffres distincts ou non.
\begin{itemize}
    \item Le choix de la lettre : 3 possibilités.
    \item Le choix des 4 chiffres : C'est une liste ordonnée avec répétitions possibles de 4 chiffres parmi les 10 chiffres (0 à 9). Le nombre de choix est donc $10^4 = 10\,000$.
    \item Par le principe multiplicatif, le nombre total de codes est : $3 \times 10^4 = 30\,000$.
\end{itemize}
\end{exemple}

\section{Exercices d'Entraînement Résolus}

\subsection*{Exercice 1 : Codes d'accès et listes}
Un code d'accès de sécurité à un bâtiment est composé de 2 lettres distinctes de l'alphabet français (parmi 26) suivies de 4 chiffres distincts ou non (de 0 à 9).
\begin{enumerate}
    \item Combien de codes différents peut-on former au total ?
    \item Combien de ces codes commencent par la lettre A et se terminent par un chiffre impair ?
\end{enumerate}

\textbf{Solution :}
\begin{enumerate}
    \item Le choix des 2 lettres distinctes ordonnées correspond à un arrangement de 2 éléments parmi 26, soit :
    \[
    A_{26}^2 = 26 \times 25 = 650
    \]
    Le choix des 4 chiffres correspond à une liste de longueur 4 d'éléments choisis parmi les 10 chiffres possibles (de 0 à 9), soit :
    \[
    10^4 = 10\,000
    \]
    Par le principe multiplicatif, le nombre total de codes possibles est :
    \[
    N = A_{26}^2 \times 10^4 = 650 \times 10\,000 = 6\,500\,000
    \]
    
    \item Si le code commence par la lettre A, il n'y a qu'une seule possibilité pour la première lettre. La deuxième lettre doit être choisie parmi les 25 lettres restantes. Le choix des lettres offre donc $1 \times 25 = 25$ possibilités.
    
    Pour les chiffres :
    - Les 3 premiers chiffres peuvent être quelconques (de 0 à 9), ce qui donne $10^3 = 1000$ possibilités.
    - Le dernier chiffre doit être impair (1, 3, 5, 7 ou 9), ce qui offre 5 possibilités.
    
    Par le principe multiplicatif, le nombre de ces codes est :
    \[
    N' = 25 \times 10^3 \times 5 = 125\,000
    \]
\end{enumerate}

\subsection*{Exercice 2 : Ordre et permutations}
Lors d'une finale de course de 100 mètres, 8 athlètes participent. On suppose qu'il n'y a pas d'ex-æquo à l'arrivée.
\begin{enumerate}
    \item Combien de podiums différents (les 3 premières places : Or, Argent, Bronze) sont possibles ?
    \item De combien de façons différentes les 8 athlètes peuvent-ils franchir la ligne d'arrivée ?
\end{enumerate}

\textbf{Solution :}
\begin{enumerate}
    \item Un podium correspond à un choix ordonné de 3 athlètes distincts parmi les 8 participants. Il s'agit d'un arrangement de 3 éléments parmi 8, soit :
    \[
    A_8^3 = 8 \times 7 \times 6 = 336 \text{ podiums possibles}
    \]
    
    \item L'ordre de franchissement de la ligne d'arrivée pour les 8 athlètes est une permutation de l'ensemble des 8 athlètes. Le nombre de façons différentes est donc :
    \[
    8! = 8 \times 7 \times 6 \times 5 \times 4 \times 3 \times 2 \times 1 = 40\,320 \text{ ordres possibles}
    \]
\end{enumerate}

\subsection*{Exercice 3 : Calculs de coefficients binomiaux et propriétés}
\begin{enumerate}
    \item Calculer sans calculatrice $\binom{10}{3}$ et $\binom{12}{10}$.
    \item Démontrer que pour tous entiers $n$ et $k$ tels que $1 \le k \le n$, on a la relation :
    \[
    k \binom{n}{k} = n \binom{n-1}{k-1}
    \]
\end{enumerate}

\textbf{Solution :}
\begin{enumerate}
    \item En appliquant la définition des coefficients binomiaux :
    \[
    \binom{10}{3} = \frac{10 \times 9 \times 8}{3 \times 2 \times 1} = \frac{720}{6} = 120
    \]
    Par la relation de symétrie, $\binom{12}{10} = \binom{12}{12-10} = \binom{12}{2}$ :
    \[
    \binom{12}{10} = \binom{12}{2} = \frac{12 \times 11}{2 \times 1} = 6 \times 11 = 66
    \]
    
    \item Exprimons le membre de gauche à l'aide des factorielles :
    \[
    k \binom{n}{k} = k \times \frac{n!}{k!(n-k)!}
    \]
    Comme $k! = k \times (k-1)!$, on peut simplifier par $k$ (puisque $k \ge 1$) :
    \[
    k \binom{n}{k} = \frac{n!}{(k-1)!(n-k)!}
    \]
    Exprimons maintenant le membre de droite :
    \[
    n \binom{n-1}{k-1} = n \times \frac{(n-1)!}{(k-1)!((n-1)-(k-1))!} = \frac{n \times (n-1)!}{(k-1)!(n-k)!} = \frac{n!}{(k-1)!(n-k)!}
    \]
    On constate que les deux expressions sont identiques, ce qui prouve l'égalité.
\end{enumerate}

\subsection*{Exercice 4 : Comité sous contraintes}
Un club de sport compte 15 membres, composé de 9 hommes et 6 femmes. On souhaite élire un bureau de 4 personnes de rôles identiques.
\begin{enumerate}
    \item Combien de bureaux différents peut-on former ?
    \item Combien de bureaux comportent exactement 2 hommes et 2 femmes ?
    \item Combien de bureaux comportent au moins une femme ?
\end{enumerate}

\textbf{Solution :}
\begin{enumerate}
    \item L'ordre des membres n'ayant pas d'importance, un bureau de 4 personnes choisi parmi les 15 membres est une combinaison de 4 éléments parmi 15. Le nombre de bureaux est :
    \[
    \binom{15}{4} = \frac{15 \times 14 \times 13 \times 12}{4 \times 3 \times 2 \times 1} = 1365
    \]
    
    \item Pour élire un bureau composé de 2 hommes et 2 femmes, on choisit simultanément 2 hommes parmi les 9 hommes et 2 femmes parmi les 6 femmes. Le nombre de façons est :
    \[
    \binom{9}{2} \times \binom{6}{2} = \frac{9 \times 8}{2} \times \frac{6 \times 5}{2} = 36 \times 15 = 540
    \]
    
    \item Calculons par le complémentaire. Les bureaux ne contenant aucune femme sont composés uniquement d'hommes. Le nombre de tels bureaux est le choix de 4 hommes parmi les 9 disponibles :
    \[
    \binom{9}{4} = \frac{9 \times 8 \times 7 \times 6}{4 \times 3 \times 2 \times 1} = 126
    \]
    Le nombre de bureaux contenant au moins une femme est donc :
    \[
    \text{Nombre total} - \text{Nombre de bureaux sans femme} = 1365 - 126 = 1239
    \]
\end{enumerate}

\subsection*{Exercice 5 : Résolution d'équations avec coefficients binomiaux}
Résoudre dans $\N$ les équations suivantes :
\begin{enumerate}
    \item $\binom{n}{2} = 45$ (avec $n \ge 2$)
    \item $\binom{n}{1} + \binom{n}{2} = 2n$ (avec $n \ge 2$)
\end{enumerate}

\textbf{Solution :}
\begin{enumerate}
    \item Pour $n \ge 2$, l'équation s'écrit :
    \[
    \frac{n(n-1)}{2} = 45 \iff n^2 - n = 90 \iff n^2 - n - 90 = 0
    \]
    C'est une équation du second degré. Calculons le discriminant :
    \[
    \Delta = (-1)^2 - 4 \times 1 \times (-90) = 1 + 360 = 361 = 19^2
    \]
    Les racines sont :
    \[
    n_1 = \frac{1 - 19}{2} = -9 \quad \text{et} \quad n_2 = \frac{1 + 19}{2} = 10
    \]
    Comme $n \in \N$, la seule solution possible est $n=10$.
    
    \item Pour $n \ge 2$, on a $\binom{n}{1} = n$ et $\binom{n}{2} = \frac{n(n-1)}{2}$. L'équation devient :
    \[
    n + \frac{n(n-1)}{2} = 2n \iff \frac{n(n-1)}{2} = n
    \]
    Comme $n \ge 2$, $n$ est non nul. On peut diviser les deux membres par $n$ :
    \[
    \frac{n-1}{2} = 1 \iff n-1 = 2 \iff n = 3
    \]
    L'unique solution est $n=3$.
\end{enumerate}

\subsection*{Exercice 6 : Dénombrement dans un jeu de cartes}
On tire simultanément 5 cartes d'un jeu de 32 cartes (qui comprend 4 couleurs : Cœur, Carreau, Pique, Trèfle, avec 8 cartes par couleur).
\begin{enumerate}
    \item Combien de tirages différents (mains) y a-t-il au total ?
    \item Combien de mains contiennent uniquement des cartes de Cœur ?
    \item Combien de mains contiennent exactement 3 Piques et 2 Trèfles ?
\end{enumerate}

\textbf{Solution :}
\begin{enumerate}
    \item Un tirage simultané de 5 cartes correspond à une combinaison de 5 cartes parmi les 32 cartes du jeu. Le nombre de tirages est :
    \[
    \binom{32}{5} = \frac{32 \times 31 \times 30 \times 29 \times 28}{5 \times 4 \times 3 \times 2 \times 1} = 201\,376
    \]
    
    \item Une main contenant uniquement des Cœurs est obtenue en choisissant 5 cartes parmi les 8 cartes de Cœur disponibles :
    \[
    \binom{8}{5} = \binom{8}{3} = \frac{8 \times 7 \times 6}{3 \times 2 \times 1} = 56
    \]
    
    \item Pour obtenir exactement 3 Piques et 2 Trèfles, on choisit simultanément 3 Piques parmi les 8 Piques possibles et 2 Trèfles parmi les 8 Trèfles possibles. Le nombre de mains est :
    \[
    \binom{8}{3} \times \binom{8}{2} = 56 \times \frac{8 \times 7}{2} = 56 \times 28 = 1568
    \]
\end{enumerate}

\subsection*{Exercice 7 : Démonstration de la formule de Pascal}
Démontrer par le calcul la formule de Pascal :
\[
\binom{n}{p} + \binom{n}{p+1} = \binom{n+1}{p+1}
\]
pour tout $n \in \N$ et $p \in \N$ tel que $0 \le p \le n-1$.

\textbf{Solution :}
Exprimons la somme du membre de gauche à l'aide des factorielles :
\[
\binom{n}{p} + \binom{n}{p+1} = \frac{n!}{p!(n-p)!} + \frac{n!}{(p+1)!(n-p-1)!}
\]
Mettons ces deux fractions au même dénominateur. On remarque que :
- $(p+1)! = (p+1) \times p!$
- $(n-p)! = (n-p) \times (n-p-1)!$
Le dénominateur commun est donc $(p+1)!(n-p)!$. Multiplions la première fraction par $\frac{p+1}{p+1}$ et la deuxième par $\frac{n-p}{n-p}$ :
\[
\begin{aligned}
\binom{n}{p} + \binom{n}{p+1} &= \frac{n!(p+1)}{(p+1)p!(n-p)!} + \frac{n!(n-p)}{(p+1)!(n-p)(n-p-1)!} \\
&= \frac{n!(p+1)}{(p+1)!(n-p)!} + \frac{n!(n-p)}{(p+1)!(n-p)!} \\
&= \frac{n!\big( (p+1) + (n-p) \big)}{(p+1)!(n-p)!} \\
&= \frac{n!(n+1)}{(p+1)!(n-p)!} \\
&= \frac{(n+1)!}{(p+1)!\big( (n+1) - (p+1) \big)!} \\
&= \binom{n+1}{p+1}
\end{aligned}
\]
La formule est ainsi démontrée.

\subsection*{Exercice 8 : Anagrammes avec répétitions de lettres}
Un anagramme d'un mot est un mot formé des mêmes lettres (ayant un sens ou non).
\begin{enumerate}
    \item Déterminer le nombre d'anagrammes du mot MATHS.
    \item Déterminer le nombre d'anagrammes du mot ANANAS.
\end{enumerate}

\textbf{Solution :}
\begin{enumerate}
    \item Le mot MATHS comporte 5 lettres toutes distinctes (M, A, T, H, S). Un anagramme correspond à une permutation de ces 5 lettres. Le nombre d'anagrammes est :
    \[
    5! = 120
    \]
    
    \item Le mot ANANAS comporte 6 lettres au total. Cependant, certaines lettres sont répétées :
    - La lettre A apparaît 3 fois.
    - La lettre N apparaît 2 fois.
    - La lettre S apparaît 1 fois.
    
    Si toutes les lettres étaient distinctes, nous aurions $6!$ permutations. Pour obtenir les anagrammes uniques, nous devons diviser par le nombre de permutations possibles des lettres répétées (qui ne changent pas le mot) :
    \[
    N = \frac{6!}{3! \times 2! \times 1!} = \frac{720}{6 \times 2 \times 1} = \frac{720}{12} = 60
    \]
    Il y a 60 anagrammes possibles pour le mot ANANAS.
\end{enumerate}

\subsection*{Exercice 9 : Nombre de sous-ensembles (Somme des coefficients binomiaux)}
Soit $E$ un ensemble fini à $n$ éléments.
\begin{enumerate}
    \item Justifier que le nombre de sous-ensembles de $E$ contenant exactement $k$ éléments ($0 \le k \le n$) est $\binom{n}{k}$.
    \item En utilisant la formule du binôme de Newton, en déduire le nombre total de sous-ensembles de $E$, c'est-à-dire la valeur de la somme :
    \[
    S = \sum_{k=0}^n \binom{n}{k}
    \]
\end{enumerate}

\textbf{Solution :}
\begin{enumerate}
    \item Choisir un sous-ensemble (ou partie) à $k$ éléments dans un ensemble à $n$ éléments revient à choisir $k$ éléments parmi $n$ sans ordre et sans répétition. C'est exactement le nombre de combinaisons de $k$ éléments parmi $n$, soit $\binom{n}{k}$.
    
    \item Le nombre total de sous-ensembles de $E$ est la somme du nombre de sous-ensembles à 0 élément, à 1 élément, ..., à $n$ éléments :
    \[
    S = \sum_{k=0}^n \binom{n}{k}
    \]
    Rappelons la formule du binôme de Newton pour tous réels $a$ et $b$ :
    \[
    (a+b)^n = \sum_{k=0}^n \binom{n}{k} a^k b^{n-k}
    \]
    En appliquant cette formule avec $a = 1$ et $b = 1$, nous obtenons :
    \[
    (1+1)^n = \sum_{k=0}^n \binom{n}{k} 1^k 1^{n-k} = \sum_{k=0}^n \binom{n}{k}
    \]
    D'où :
    \[
    S = 2^n
    \]
    Un ensemble à $n$ éléments possède donc $2^n$ sous-ensembles au total.
\end{enumerate}

\subsection*{Exercice 10 : Tirage d'urnes sous conditions}
Une urne contient 12 boules indiscernables au toucher : 5 rouges, 4 vertes et 3 jaunes. On tire simultanément 3 boules de l'urne.
\begin{enumerate}
    \item Combien de tirages possibles y a-t-il au total ?
    \item Combien de tirages contiennent exactement une boule de chaque couleur ?
    \item Combien de tirages ne contiennent aucune boule rouge ?
\end{enumerate}

\textbf{Solution :}
\begin{enumerate}
    \item Le tirage est simultané, l'ordre n'importe pas. Un tirage est une combinaison de 3 boules parmi les 12 boules de l'urne :
    \[
    \binom{12}{3} = \frac{12 \times 11 \times 10}{3 \times 2 \times 1} = 2 \times 11 \times 10 = 220 \text{ tirages possibles}
    \]
    
    \item Un tirage contenant une boule de chaque couleur correspond au choix d'une rouge parmi 5, une verte parmi 4 et une jaune parmi 3 :
    \[
    \binom{5}{1} \times \binom{4}{1} \times \binom{3}{1} = 5 \times 4 \times 3 = 60 \text{ tirages}
    \]
    
    \item Ne contenir aucune boule rouge signifie que les 3 boules sont choisies parmi les boules vertes ou jaunes (soit $4 + 3 = 7$ boules non rouges) :
    \[
    \binom{7}{3} = \frac{7 \times 6 \times 5}{3 \times 2 \times 1} = 35 \text{ tirages}
    \]
\end{enumerate}

\section{Problèmes Résolus}

\subsection*{Problème 1 : Dénombrement des mains au poker}
Dans un jeu standard de 52 cartes (comprenant 4 couleurs de 13 cartes chacune), une main est formée de 5 cartes distribuées simultanément.
\begin{enumerate}
    \item Calculer le nombre total de mains de 5 cartes possibles.
    \item Déterminer le nombre de mains contenant les combinaisons suivantes :
    \begin{enumerate}
        \item \textbf{Une paire} : exactement deux cartes de même valeur, et trois autres cartes de valeurs distinctes entre elles et différentes de la paire.
        \item \textbf{Un brelan} : exactement trois cartes de même valeur, et deux autres de valeurs distinctes.
        \item \textbf{Une couleur} : cinq cartes de la même couleur (parmi Cœur, Carreau, Pique, Trèfle).
        \item \textbf{Un full} : trois cartes de même valeur (un brelan) et deux autres d'une autre même valeur (une paire).
        \item \textbf{Un carré} : quatre cartes de même valeur.
    \end{enumerate}
\end{enumerate}

\textbf{Solution :}
\begin{enumerate}
    \item Une main est une combinaison de 5 cartes parmi 52. Le nombre de mains possibles est :
    \[
    \binom{52}{5} = \frac{52 \times 51 \times 30 \times 29 \times 28}{5 \times 4 \times 3 \times 2 \times 1} \text{ (après simplification)} = 2\,598\,960
    \]
    
    \item Déterminons le nombre de mains pour chaque configuration :
    \begin{enumerate}
        \item \textbf{Une paire} :
        - On choisit la valeur de la paire parmi les 13 valeurs possibles : $\binom{13}{1} = 13$ choix.
        - On choisit les 2 cartes de cette valeur parmi les 4 enseignes possibles : $\binom{4}{2} = 6$ choix.
        - On choisit les 3 autres valeurs distinctes parmi les 12 valeurs restantes : $\binom{12}{3} = 220$ choix.
        - Pour chacune de ces 3 valeurs, on choisit une carte parmi les 4 enseignes possibles : $\binom{4}{1}^3 = 4^3 = 64$ choix.
        Par le principe multiplicatif, le nombre de mains est :
        \[
        N_{\text{paire}} = 13 \times 6 \times 220 \times 64 = 1\,098\,240
        \]
        
        \item \textbf{Un brelan} :
        - On choisit la valeur du brelan parmi les 13 valeurs : $\binom{13}{1} = 13$ choix.
        - On choisit les 3 cartes de cette valeur parmi les 4 enseignes : $\binom{4}{3} = 4$ choix.
        - On choisit les 2 autres valeurs distinctes parmi les 12 restantes : $\binom{12}{2} = 66$ choix.
        - Pour chacune de ces 2 valeurs, on choisit une carte parmi les 4 enseignes : $\binom{4}{1}^2 = 4^2 = 16$ choix.
        Le nombre de mains est :
        \[
        N_{\text{brelan}} = 13 \times 4 \times 66 \times 16 = 54\,912
        \]
        
        \item \textbf{Une couleur} :
        - On choisit l'une des 4 couleurs (Cœur, Carreau, Pique, Trèfle) : $\binom{4}{1} = 4$ choix.
        - On choisit les 5 cartes de cette couleur parmi les 13 cartes disponibles : $\binom{13}{5} = \frac{13 \times 12 \times 11 \times 10 \times 9}{5 \times 4 \times 3 \times 2 \times 1} = 1287$ choix.
        Le nombre de mains est :
        \[
        N_{\text{couleur}} = 4 \times 1287 = 5148
        \]
        
        \item \textbf{Un full} :
        - On choisit la valeur du brelan parmi les 13 valeurs : $\binom{13}{1} = 13$ choix.
        - On choisit 3 cartes de cette valeur parmi les 4 enseignes : $\binom{4}{3} = 4$ choix.
        - On choisit la valeur de la paire parmi les 12 valeurs restantes : $\binom{12}{1} = 12$ choix.
        - On choisit 2 cartes de cette valeur parmi les 4 enseignes : $\binom{4}{2} = 6$ choix.
        Le nombre de mains est :
        \[
        N_{\text{full}} = 13 \times 4 \times 12 \times 6 = 3744
        \]
        
        \item \textbf{Un carré} :
        - On choisit la valeur du carré parmi les 13 valeurs : $\binom{13}{1} = 13$ choix.
        - On prend les 4 cartes de cette valeur : $\binom{4}{4} = 1$ choix.
        - On choisit la 5-ème carte restante parmi les 48 cartes restantes de valeurs différentes : $\binom{48}{1} = 48$ choix.
        Le nombre de mains est :
        \[
        N_{\text{carré}} = 13 \times 1 \times 48 = 624
        \]
    \end{enumerate}
\end{enumerate}

\subsection*{Problème 2 : Chemins sur un réseau quadrillé}
On considère un réseau quadrillé du plan. Un chemin minimal reliant le point $O(0,0)$ au point $M(a,b)$, où $a$ et $b$ sont des entiers naturels, est constitué d'une succession de pas unitaires : soit vers la droite (noté D, incrémentant l'abscisse de 1), soit vers le haut (noté H, incrémentant l'ordonnée de 1).
\begin{enumerate}
    \item Justifier qu'un tel chemin comporte exactement $a+b$ pas, dont $a$ pas vers la droite. En déduire que le nombre total de chemins possibles est $\binom{a+b}{a}$.
    \item On pose $a=5$ et $b=4$. Calculer le nombre de chemins reliant $(0,0)$ à $(5,4)$.
    \item On impose de passer par le point d'étape $A(2,2)$. Combien de chemins reliant $(0,0)$ à $(5,4)$ passent par $A$ ?
    \item Combien de chemins ne passent pas par $A$ ?
    \item Combien de chemins reliant $(0,0)$ à $(5,4)$ passent par le point $A(2,2)$ ou par le point $B(3,3)$ ?
\end{enumerate}

\textbf{Solution :}
\begin{enumerate}
    \item Pour se rendre de $O(0,0)$ à $M(a,b)$ par des déplacements minimaux, on doit nécessairement effectuer exactement $a$ pas vers la droite (pour atteindre l'abscisse $a$) et $b$ pas vers le haut (pour atteindre l'ordonnée $b$).
    Le nombre total de pas est donc $a+b$. Un chemin est entièrement caractérisé par l'emplacement des $a$ pas vers la droite parmi les $a+b$ pas totaux. Le nombre de façons de choisir ces emplacements est donné par la combinaison :
    \[
    \binom{a+b}{a}
    \]
    (qui est égal par symétrie à $\binom{a+b}{b}$, c'est-à-dire le choix des pas vers le haut).
    
    \item Pour $a=5$ et $b=4$, le nombre de chemins possibles est :
    \[
    \binom{5+4}{5} = \binom{9}{5} = \frac{9 \times 8 \times 7 \times 6 \times 5}{5 \times 4 \times 3 \times 2 \times 1} = 126
    \]
    
    \item Un chemin passant par $A(2,2)$ est composé :
    - d'un chemin de $O(0,0)$ à $A(2,2)$ : ici $a_1=2$ et $b_1=2$, ce qui donne $\binom{2+2}{2} = \binom{4}{2} = 6$ chemins.
    - d'un chemin de $A(2,2)$ à $M(5,4)$ : le déplacement requis est de $5-2=3$ pas vers la droite et $4-2=2$ pas vers le haut. Il y a donc $\binom{3+2}{3} = \binom{5}{3} = 10$ chemins.
    
    Par le principe multiplicatif, le nombre de chemins de $O$ à $M$ passant par $A$ est :
    \[
    N_A = \binom{4}{2} \times \binom{5}{3} = 6 \times 10 = 60
    \]
    
    \item Le nombre de chemins ne passant pas par $A$ est le complémentaire du nombre de chemins passant par $A$ :
    \[
    N_{\bar{A}} = \text{Total} - N_A = 126 - 60 = 66
    \]
    
    \item Soit $\mathcal{C}_A$ l'ensemble des chemins passant par $A$ et $\mathcal{C}_B$ l'ensemble des chemins passant par $B$. On cherche $\text{Card}(\mathcal{C}_A \cup \mathcal{C}_B)$.
    D'après le principe additif :
    \[
    \text{Card}(\mathcal{C}_A \cup \mathcal{C}_B) = \text{Card}(\mathcal{C}_A) + \text{Card}(\mathcal{C}_B) - \text{Card}(\mathcal{C}_A \cap \mathcal{C}_B)
    \]
    - On a déjà $\text{Card}(\mathcal{C}_A) = 60$.
    - Calculons $\text{Card}(\mathcal{C}_B)$ :
      * Chemin de $O(0,0)$ à $B(3,3)$ : $\binom{3+3}{3} = \binom{6}{3} = 20$.
      * Chemin de $B(3,3)$ à $M(5,4)$ : déplacement de $5-3=2$ vers la droite et $4-3=1$ vers le haut, soit $\binom{2+1}{2} = \binom{3}{2} = 3$.
      * D'où $\text{Card}(\mathcal{C}_B) = 20 \times 3 = 60$.
    - Calculons $\text{Card}(\mathcal{C}_A \cap \mathcal{C}_B)$, c'est-à-dire les chemins passant par $A$ ET par $B$. Comme $A(2,2)$ est situé avant $B(3,3)$ ($2 \le 3$ en abscisse et ordonnée), un tel chemin va de $O$ à $A$, puis de $A$ à $B$, puis de $B$ à $M$ :
      * De $O(0,0)$ à $A(2,2)$ : $\binom{4}{2} = 6$ chemins.
      * De $A(2,2)$ à $B(3,3)$ : déplacement de $3-2=1$ vers la droite et $3-2=1$ vers le haut, soit $\binom{1+1}{1} = \binom{2}{1} = 2$ chemins.
      * De $B(3,3)$ à $M(5,4)$ : $\binom{3}{2} = 3$ chemins.
      * Ainsi, $\text{Card}(\mathcal{C}_A \cap \mathcal{C}_B) = 6 \times 2 \times 3 = 36$.
      
    On en déduit :
    \[
    \text{Card}(\mathcal{C}_A \cup \mathcal{C}_B) = 60 + 60 - 36 = 84
    \]
    Il y a 84 chemins qui passent par $A$ ou par $B$.
\end{enumerate}

\subsection*{Problème 3 : Identités binomiales combinatoires et algébriques}
Soit $n$ un entier naturel non nul. Le but de ce problème est d'étudier la somme :
\[
S_n = \sum_{k=1}^n k \binom{n}{k}
\]
\begin{enumerate}
    \item Calculer $S_1$, $S_2$ et $S_3$. Conjecturer une formule simplifiée pour $S_n$.
    \item \textbf{Méthode algébrique} : En utilisant la relation d'égalité de l'exercice 3 ($k \binom{n}{k} = n \binom{n-1}{k-1}$), démontrer la conjecture.
    \item \textbf{Méthode combinatoire} : On considère un groupe de $n$ personnes. On souhaite former un comité de taille quelconque (contenant au moins une personne) et choisir un président au sein de ce comité.
    \begin{enumerate}
        \item En dénombrant selon la taille $k$ du comité choisi, montrer que le nombre total de choix possibles est égal à $S_n$.
        \item En choisissant d'abord le président, puis le reste du comité parmi les personnes restantes, déterminer à nouveau le nombre total de choix possibles.
        \item Conclure.
    \end{enumerate}
\end{enumerate}

\textbf{Solution :}
\begin{enumerate}
    \item Calculons les premières valeurs :
    - Pour $n=1$ : $S_1 = 1 \binom{1}{1} = 1 \times 1 = 1$.
    - Pour $n=2$ : $S_2 = 1 \binom{2}{1} + 2 \binom{2}{2} = 1 \times 2 + 2 \times 1 = 4$.
    - Pour $n=3$ : $S_3 = 1 \binom{3}{1} + 2 \binom{3}{2} + 3 \binom{3}{3} = 1 \times 3 + 2 \times 3 + 3 \times 1 = 12$.
    
    On remarque que :
    - $S_1 = 1 = 1 \times 2^0$
    - $S_2 = 4 = 2 \times 2^1$
    - $S_3 = 12 = 3 \times 2^2$
    
    On conjecture ainsi la formule générale :
    \[
    S_n = n 2^{n-1}
    \]
    
    \item D'après la relation $k \binom{n}{k} = n \binom{n-1}{k-1}$, on a :
    \[
    S_n = \sum_{k=1}^n k \binom{n}{k} = \sum_{k=1}^n n \binom{n-1}{k-1} = n \sum_{k=1}^n \binom{n-1}{k-1}
    \]
    Effectuons le changement d'indice $j = k-1$. Quand $k$ varie de $1$ à $n$, l'indice $j$ varie de $0$ à $n-1$ :
    \[
    S_n = n \sum_{j=0}^{n-1} \binom{n-1}{j}
    \]
    Or, d'après l'exercice 9, la somme des coefficients binomiaux d'ordre $n-1$ vaut $\sum_{j=0}^{n-1} \binom{n-1}{j} = 2^{n-1}$.
    On en déduit immédiatement :
    \[
    S_n = n 2^{n-1}
    \]
    
    \item Utilisons le raisonnement par double dénombrement :
    \begin{enumerate}
        \item Pour former un comité avec un président, on peut d'abord choisir un comité de $k$ personnes parmi les $n$ ($\binom{n}{k}$ choix), puis désigner un président parmi les $k$ membres élus ($k$ choix). Le nombre de choix pour un comité de taille $k$ est $k\binom{n}{k}$.
        Comme la taille $k$ du comité peut être n'importe quelle valeur de $1$ à $n$, le nombre total de choix possibles est :
        \[
        \sum_{k=1}^n k \binom{n}{k} = S_n
        \]
        
        \item On peut procéder autrement :
        - On choisit d'abord le président du comité parmi les $n$ personnes du groupe ($n$ choix).
        - Ensuite, pour chacune des $n-1$ personnes restantes, on décide si elle fait partie ou non du reste du comité (2 possibilités par personne : soit elle y est, soit elle n'y est pas). Cela correspond à former un sous-ensemble quelconque de l'ensemble des $n-1$ personnes restantes. Le nombre de choix est donc $2^{n-1}$.
        Par le principe multiplicatif, le nombre total de choix possibles est :
        \[
        n \times 2^{n-1}
        \]
        
        \item Les deux méthodes dénombrant le même ensemble de configurations, on en déduit l'égalité :
        \[
        S_n = n 2^{n-1}
        \]
    \end{enumerate}
\end{enumerate}

\subsection*{Problème 4 : Distribution d'objets dans des boîtes (Partitions et Étoiles et barres)}
On souhaite étudier la distribution de $n$ objets dans $k$ boîtes distinctes (numérotées de 1 à $k$).
\begin{enumerate}
    \item On suppose d'abord que les $n$ objets sont **distincts** (par exemple, des lettres différentes à poster).
    Justifier que le nombre total de répartitions possibles est $k^n$.
    
    \item On suppose désormais que les $n$ objets sont **identiques** (par exemple, des bonbons identiques). On cherche à déterminer le nombre de répartitions.
    \begin{enumerate}
        \item Dans le cas $n=4$ et $k=3$, on note $(x_1, x_2, x_3)$ la répartition où $x_i$ est le nombre d'objets dans la boîte $i$ ($x_i \in \N$ et $x_1+x_2+x_3=4$).
        Lister toutes les répartitions possibles (les triplets) et en donner le nombre.
        
        \item De façon générale, on représente une répartition en plaçant les $n$ objets sous la forme de ronds $\bullet$ séparés par $k-1$ barres verticales $|$ qui délimitent les $k$ boîtes.
        Par exemple, pour $n=4, k=3$, la répartition $(2, 0, 2)$ est représentée par :
        \[
        \bullet \bullet \mid \mid \bullet \bullet
        \]
        Justifier que le nombre total de symboles (ronds et barres) dans cette représentation est $n+k-1$.
        
        \item Expliquer pourquoi dénombrer les répartitions possibles revient à choisir les positions des $n$ ronds parmi les $n+k-1$ emplacements possibles. En déduire que le nombre de répartitions est égal à :
        \[
        \binom{n+k-1}{n}
        \]
        
        \item Calculer le nombre de façons de distribuer 10 bonbons identiques à 3 enfants.
    \end{enumerate}
\end{enumerate}

\textbf{Solution :}
\begin{enumerate}
    \item Pour distribuer $n$ objets distincts dans $k$ boîtes distinctes, chaque objet peut être placé dans n'importe laquelle des $k$ boîtes. Il y a donc $k$ choix pour le premier objet, $k$ choix pour le second, ..., et $k$ choix pour le $n$-ième.
    Par le principe multiplicatif, le nombre de répartitions possibles est :
    \[
    \underbrace{k \times k \times \dots \times k}_{n \text{ fois}} = k^n
    \]
    
    \item Distribution d'objets identiques :
    \begin{enumerate}
        \item Les répartitions possibles pour $n=4$ objets dans $k=3$ boîtes sont les triplets d'entiers naturels $(x_1, x_2, x_3)$ vérifiant $x_1+x_2+x_3=4$ :
        - Les permutations de $(4,0,0)$ : $(4,0,0)$, $(0,4,0)$, $(0,0,4)$ (3 triplets) ;
        - Les permutations de $(3,1,0)$ : $(3,1,0)$, $(3,0,1)$, $(1,3,0)$, $(1,0,3)$, $(0,3,1)$, $(0,1,3)$ (6 triplets) ;
        - Les permutations de $(2,2,0)$ : $(2,2,0)$, $(2,0,2)$, $(0,2,2)$ (3 triplets) ;
        - Les permutations de $(2,1,1)$ : $(2,1,1)$, $(1,2,1)$, $(1,1,2)$ (3 triplets).
        Le nombre total de répartitions possibles est de $3 + 6 + 3 + 3 = 15$.
        
        \item Pour délimiter $k$ boîtes distinctes, il suffit d'utiliser $k-1$ séparateurs (les barres $|$). La répartition est modélisée par une suite contenant $n$ symboles $\bullet$ (pour les objets) et $k-1$ symboles $\mid$ (pour les séparateurs).
        Le nombre total de symboles est donc la somme :
        \[
        n + (k-1) = n+k-1
        \]
        
        \item Dégager une configuration unique de répartition revient à spécifier où se trouvent les $n$ symboles $\bullet$ parmi les $n+k-1$ emplacements de la suite (les emplacements non choisis contenant nécessairement des barres $\mid$).
        Le nombre de choix est donné par le nombre de combinaisons de $n$ éléments parmi $n+k-1$, soit :
        \[
        \binom{n+k-1}{n}
        \]
        Pour $n=4, k=3$, cela donne $\binom{4+3-1}{4} = \binom{6}{4} = \binom{6}{2} = 15$, ce qui corrobore le résultat de la question (a).
        
        \item Distribuer 10 bonbons identiques ($n=10$) à 3 enfants ($k=3$, assimilés à 3 boîtes distinctes) correspond au nombre de répartitions :
        \[
        \binom{10+3-1}{10} = \binom{12}{10} = \binom{12}{2} = \frac{12 \times 11}{2 \times 1} = 66
        \]
        Il y a 66 répartitions possibles.
    \end{enumerate}
\end{enumerate}

\subsection*{Problème 5 : Le problème des dérangements (du vestiaire)}
Quatre personnes déposent leur chapeau au vestiaire d'un restaurant. En repartant, le responsable leur remet au hasard un chapeau parmi les quatre.
On s'intéresse au nombre de façons de distribuer les quatre chapeaux de telle sorte que **personne** ne reçoive son propre chapeau. Une telle configuration s'appelle un dérangement de l'ensemble $\{1, 2, 3, 4\}$.
\begin{enumerate}
    \item Justifier que le nombre total de façons de distribuer les 4 chapeaux aux 4 personnes est $4!$.
    \item Pour tout entier $i \in \{1, 2, 3, 4\}$, on note $A_i$ l'ensemble des distributions dans lesquelles la personne $i$ récupère son propre chapeau.
    \begin{enumerate}
        \item Déterminer $\text{Card}(A_1)$.
        \item Déterminer $\text{Card}(A_1 \cap A_2)$.
        \item De manière générale, pour $k$ personnes fixées ($1 \le k \le 4$), quel est le cardinal de l'intersection de leurs ensembles $A_i$ correspondants ?
    \end{enumerate}
    \item On souhaite utiliser le principe d'inclusion-exclusion pour calculer le nombre de distributions dans lesquelles au moins une personne récupère son chapeau, soit $\text{Card}(A_1 \cup A_2 \cup A_3 \cup A_4)$. La formule pour 4 ensembles est :
    \[
    \begin{aligned}
    \text{Card}(A_1 \cup A_2 \cup A_3 \cup A_4) &= \sum_{i=1}^4 \text{Card}(A_i) - \sum_{1 \le i < j \le 4} \text{Card}(A_i \cap A_j) \\
    &\quad + \sum_{1 \le i < j < k \le 4} \text{Card}(A_i \cap A_j \cap A_k) - \text{Card}(A_1 \cap A_2 \cap A_3 \cap A_4)
    \end{aligned}
    \]
    \begin{enumerate}
        \item Calculer le nombre de termes dans la deuxième somme et dans la troisième somme.
        \item Calculer la valeur numérique de chacun des termes de la formule et montrer que :
        \[
        \text{Card}(A_1 \cup A_2 \cup A_3 \cup A_4) = 15
        \]
    \end{enumerate}
    \item En déduire le nombre de dérangements possibles pour $n=4$ personnes.
    \item \textbf{Généralisation} : On admet que pour $n$ personnes, le nombre de dérangements est donné par la formule :
    \[
    D_n = n! \sum_{k=0}^n \frac{(-1)^k}{k!}
    \]
    Vérifier que la formule donne bien la même valeur pour $n=4$.
\end{enumerate}

\textbf{Solution :}
\begin{enumerate}
    \item Distribuer 4 chapeaux distincts à 4 personnes distinctes correspond à ordonner les chapeaux, ce qui équivaut à une permutation des 4 éléments. Le nombre total de distributions est :
    \[
    4! = 24
    \]
    
    \item Calculons les cardinaux :
    \begin{enumerate}
        \item Dans $A_1$, la personne 1 reçoit son chapeau 1. Il reste à distribuer les 3 autres chapeaux aux 3 autres personnes de manière quelconque. Il y a donc $3! = 6$ façons.
        D'où $\text{Card}(A_1) = 6$. (Par symétrie, c'est aussi le cas pour $A_2, A_3, A_4$).
        
        \item Dans $A_1 \cap A_2$, la personne 1 reçoit son chapeau 1 et la personne 2 reçoit son chapeau 2. Il reste à distribuer les 2 chapeaux restants aux 2 personnes restantes.
        Il y a $2! = 2$ façons. D'où $\text{Card}(A_1 \cap A_2) = 2$.
        
        \item De manière générale, si $k$ personnes reçoivent leur propre chapeau, les positions de ces $k$ chapeaux sont fixées. Il reste alors à distribuer librement les $4-k$ chapeaux restants aux $4-k$ personnes restantes.
        Le nombre de distributions correspondantes est donc :
        \[
        (4-k)!
        \]
    \end{enumerate}
    
    \item Calculons la réunion par le principe d'inclusion-exclusion :
    \begin{enumerate}
        \item - Le nombre de termes dans la première somme est $\binom{4}{1} = 4$.
        - La deuxième somme porte sur les couples d'indices distincts $\{i, j\}$ tels que $1 \le i < j \le 4$. Le nombre de termes est le nombre de choix de 2 indices parmi 4, soit :
        \[
        \binom{4}{2} = 6
        \]
        - La troisième somme porte sur les triplets d'indices distincts $\{i, j, k\}$ tels que $1 \le i < j < k \le 4$. Le nombre de termes est :
        \[
        \binom{4}{3} = 4
        \]
        
        \item Évaluons chaque somme :
        - $\sum_{i=1}^4 \text{Card}(A_i) = \binom{4}{1} \times 3! = 4 \times 6 = 24$.
        - $\sum_{1 \le i < j \le 4} \text{Card}(A_i \cap A_j) = \binom{4}{2} \times 2! = 6 \times 2 = 12$.
        - $\sum_{1 \le i < j < k \le 4} \text{Card}(A_i \cap A_j \cap A_k) = \binom{4}{3} \times 1! = 4 \times 1 = 4$.
        - $\text{Card}(A_1 \cap A_2 \cap A_3 \cap A_4) = \binom{4}{4} \times 0! = 1 \times 1 = 1$.
        
        En injectant ces valeurs dans la formule :
        \[
        \text{Card}(A_1 \cup A_2 \cup A_3 \cup A_4) = 24 - 12 + 4 - 1 = 15
        \]
    \end{enumerate}
    
    \item Un dérangement est une distribution où personne ne reçoit son propre chapeau. C'est l'ensemble complémentaire de l'événement « au moins une personne reçoit son propre chapeau ».
    Le nombre de dérangements est donc :
    \[
    D_4 = \text{Total} - \text{Card}(A_1 \cup A_2 \cup A_3 \cup A_4) = 24 - 15 = 9
    \]
    Il y a 9 dérangements possibles.
    
    \item Pour $n=4$, la formule donne :
    \[
    \begin{aligned}
    D_4 &= 4! \sum_{k=0}^4 \frac{(-1)^k}{k!} \\
    &= 24 \left( \frac{(-1)^0}{0!} + \frac{(-1)^1}{1!} + \frac{(-1)^2}{2!} + \frac{(-1)^3}{3!} + \frac{(-1)^4}{4!} \right) \\
    &= 24 \left( 1 - 1 + \frac{1}{2} - \frac{1}{6} + \frac{1}{24} \right) \\
    &= 24 \left( \frac{12 - 4 + 1}{24} \right) = 9
    \end{aligned}
    \]
    La formule donne bien 9, ce qui confirme son exactitude.
\end{enumerate}
